\documentclass[10pt,a4paper]{amsart}
\usepackage[margin=19mm]{geometry}
\usepackage{amsmath,amssymb,mathtools}
\usepackage[scr=rsfs,cal=euler]{mathalfa}
\usepackage{xfrac}
\usepackage[unicode,hidelinks,bookmarks=false]{hyperref}
\newcommand{\ie}{i.e.\ }
\newcommand{\cf}{cf.\ }
\newcommand{\resp}{respectively\ }
\newcommand{\Subsection}{\subsection}
\providecommand{\nobreakdash}{\nobreak\hbox{-}}
\setlength{\emergencystretch}{2em}
\multlinegap0pt
\usepackage{enumerate}
\usepackage{amssymb}
\usepackage{mathtools}
\usepackage{float}
\usepackage{xcolor}
\newtheorem{lemma}{Lemma}
\title{Taylor Positivity of Ehrhart Polynomials}
\author{Feihu Liu$^{\color{blue} \dag}$ and Zihao Zhang$^{\color{blue} \S}$}
\date{Source: arXiv:2609.03327v1 (CC BY 4.0)}
\begin{document}
\maketitle

\noindent\emph{Source and licence.} Taylor Positivity of Ehrhart Polynomials. Version arXiv:2609.03327v1 (CC BY 4.0), distributed by arXiv under the Creative Commons Attribution 4.0 International licence, http://creativecommons.org/licenses/by/4.0/, as recorded in the arXiv licence metadata for this exact version and shown on the abstract page https://arxiv.org/abs/2609.03327v1. Full licence notice: \texttt{licenses/arxiv-2609.03327-CC-BY-4.0.txt}.\par\smallskip

\noindent\textit{Editorial scope.} Counted unit: the article's lemma lem:shifted-rising-comparison (source0.tex lines 967--982). The article's complete original proof of it is delivered with it (source0.tex lines 983--1078, one \texttt{proof} environment of 96 lines). No mathematics is written, repaired or re-derived: the statement and the proof are the article's own lines, in the article's own notation, and no retained line is substituted or re-flowed.  No further source-local context is needed: every cross-reference of every retained line resolves inside the two delivered spans.  Nothing was omitted from any retained span, so the package carries no omission record.  No retained line contains a citation, so the wrapper carries no bibliography and no bibliography entry.  No retained line contains a figure environment, an image-inclusion command or drawing code, so no image is delivered and the package declares no asset dependency.  Editorial formatting only, in addition: an AMS \texttt{amsart} preamble that restates the article's own declarations and macros used by the retained lines, namely the environments \texttt{lemma} on separate counters, together with the standard packages those lines need.  The retained environments print in this document as Lemma 1 in order of appearance, whereas the article prints the same items with the numbering of its own full document.  The complete native source of the article as distributed in the arXiv e-print for this exact version is bundled unchanged as \texttt{sources/209287/source0.tex}.
\begin{lemma}\label{lem:shifted-rising-comparison}
Let $d=2m\geq2$, let $1\leq\ell\leq d-1$, and set
\[C_{d,\ell}:=[x^\ell]x^{\overline d}.
\]
Then $C_{d,\ell}>0$, and the following two inequalities hold.
\begin{enumerate}
\item[(a)] For every integer $a$ with $0\leq a\leq m$, we have 
\begin{equation}\label{eq:positive-rising-shift}
[x^\ell](x+a)^{\overline d}\geq\frac{2a+1}{d-1}\,C_{d,\ell}.
\end{equation}
\item[(b)] For every integer $b$ with $1\leq b\leq m$, we have 
\begin{equation}\label{eq:negative-rising-shift}
[x^\ell](x-b)^{\overline d}\geq-\frac{2b-1}{d-1}\,C_{d,\ell}.
\end{equation}
\end{enumerate}
\end{lemma}

\begin{proof}
Because $x^{\overline d}=x\prod_{\nu=1}^{d-1}(x+\nu)$, every coefficient of positive degree is positive.  Hence
$C_{d,\ell}>0$ for $1\leq\ell\leq d$.
We first prove part~(a).  Set $q=d-\ell$.  Since $1\leq\ell\leq d-1$, we have $1\leq q\leq d-1$.  Define
$E_q(a):=e_q(a,a+1,\ldots,a+d-1)$.
As usual, $e_0=1$. Expanding the product $(x+a)^{\overline d}$ shows that
\begin{equation}\label{eq:Eq-coefficient}
 E_q(a)=[x^\ell](x+a)^{\overline d}.
\end{equation}
Moreover, $E_q(0)=e_q(0,1,\ldots,d-1)=e_q(1,\ldots,d-1)=C_{d,\ell}$.
Every summand in
\[E_q(a)=\sum_{\substack{S\subseteq\{0,\ldots,d-1\}\\|S|=q}}\prod_{\nu\in S}(a+\nu)\]
is a polynomial in $a$ with nonnegative coefficients.  Consequently $E_q(a)$ itself has nonnegative coefficients, and hence
\begin{equation}\label{eq:Eq-linear-lower-bound}
 E_q(a)\geq E_q(0)+aE_q'(0)\qquad(a\geq0).
\end{equation}

We next estimate $E_q'(0)$.  Differentiating an elementary symmetric
function after shifting all its variables gives
\[E_q'(a)=\sum_{\nu=0}^{d-1}e_{q-1}(a,a+1,\ldots,\widehat{a+\nu},\ldots,a+d-1),
\]
where the hat denotes omission.  A fixed monomial of degree $q-1$
occurs in exactly $d-(q-1)=d-q+1$ of the summands.  Therefore, we have 
\begin{equation}\label{eq:Eq-derivative}
 E_q'(a)=(d-q+1)e_{q-1}(a,a+1,\ldots,a+d-1).
\end{equation}
At $a=0$, this becomes $E_q'(0)=(d-q+1)e_{q-1}(1,\ldots,d-1)$.

We claim that
\begin{equation}\label{eq:elementary-ratio-bound}
 e_q(1,\ldots,d-1)
 \leq\frac{(d-q+1)(d-1)}2e_{q-1}(1,\ldots,d-1).
\end{equation}
To prove the claim, set $y_\nu=\nu$ for $1\leq\nu\leq d-1$.
Double counting pairs consisting of a $q$-element set and one of its
elements gives
\begin{align*}
 q e_q(y_1,\ldots,y_{d-1})
=\sum_{\substack{T\subseteq\{1,\ldots,d-1\}\\|T|=q-1}}\left(\prod_{\nu\in T}y_\nu\right)\left(\sum_{\mu\notin T}y_\mu\right)
&\leq\left(\sum_{\mu=1}^{d-1}\mu\right)e_{q-1}(y_1,\ldots,y_{d-1})
\\& =\frac{d(d-1)}2e_{q-1}(y_1,\ldots,y_{d-1}).
\end{align*}
Thus
\[e_q(1,\ldots,d-1)\leq\frac{d(d-1)}{2q}e_{q-1}(1,\ldots,d-1).
\]
Finally, we have $q(d-q+1)-d=(q-1)(d-q)\geq0$, so $d/q\leq d-q+1$.  This proves \eqref{eq:elementary-ratio-bound}.
Combining \eqref{eq:Eq-derivative} and \eqref{eq:elementary-ratio-bound}, we obtain
$E_q'(0)\geq\frac2{d-1}E_q(0)$.
Substitution into \eqref{eq:Eq-linear-lower-bound} yields
\begin{align*}
E_q(a)\geq E_q(0)\left(1+\frac{2a}{d-1}\right)=\frac{d-1+2a}{d-1}E_q(0)\geq\frac{2a+1}{d-1}E_q(0),
\end{align*}
where the last inequality uses $d\geq2$.  Recalling \eqref{eq:Eq-coefficient} and $E_q(0)=C_{d,\ell}$ proves \eqref{eq:positive-rising-shift}.

We now prove part~(b).  Fix $1\leq b\leq m=d/2$, and set $\gamma_b:=\frac{2b-1}{d-1}$.
The two rising factorials have the factorizations
\begin{align*}
 (x-b)^{\overline d}=x^{\overline{d-b}}\prod_{\nu=1}^{b}(x-\nu),\quad
 x^{\overline d}=x^{\overline{d-b}}\prod_{\nu=d-b}^{d-1}(x+\nu).
\end{align*}
Therefore, we have 
\begin{equation}\label{eq:tail-factorization}
 (x-b)^{\overline d}+\gamma_bx^{\overline d}=x^{\overline{d-b}}H_{d,b}(x),
\end{equation}
where
\[ H_{d,b}(x) :=\prod_{\nu=1}^{b}(x-\nu) +\gamma_b\prod_{\nu=d-b}^{d-1}(x+\nu).
\]
We show that every coefficient of $H_{d,b}(x)$ is nonnegative.  For $0\leq j\leq b$, the coefficient of $x^{b-j}$ is
\begin{equation}\label{eq:H-coefficient}
 (-1)^je_j(1,\ldots,b)+\gamma_be_j(d-b,\ldots,d-1).
\end{equation}
If $j$ is even, this is plainly nonnegative.  Suppose that $j$ is odd.
For $1\leq\nu\leq b$, one has
\[ b(d-b+\nu-1)-(d-1)\nu =(d-b-1)(b-\nu)\geq0.
\]
Thus
\begin{equation}\label{eq:entrywise-high-low}
d-b+\nu-1\geq\frac{d-1}{b}\nu\qquad(1\leq\nu\leq b).
\end{equation}
Multiplying \eqref{eq:entrywise-high-low} over each $j$-element subset and then summing gives
\begin{equation}\label{eq:elementary-high-low}
 e_j(d-b,\ldots,d-1)\geq\left(\frac{d-1}{b}\right)^je_j(1,\ldots,b).
\end{equation}
Since $b\leq d/2$, we have $(d-1)/b\geq1$.  Since $j\geq1$,
\[\left(\frac{d-1}{b}\right)^j\geq\frac{d-1}{b}\geq\frac{d-1}{2b-1}=\frac1{\gamma_b}.
\]
It follows from \eqref{eq:elementary-high-low} that
\[\gamma_be_j(d-b,\ldots,d-1)\geq e_j(1,\ldots,b).
\]
Hence \eqref{eq:H-coefficient} is nonnegative also when $j$ is odd.
We have proved that $H_{d,b}(x)$ has nonnegative coefficients.  The polynomial $x^{\overline{d-b}}$ also has nonnegative coefficients. Equation \eqref{eq:tail-factorization} therefore shows, coefficient by
coefficient, that
\[[x^\ell](x-b)^{\overline d}+\gamma_b[x^\ell]x^{\overline d}\geq0.
\]
Since $[x^\ell]x^{\overline d}=C_{d,\ell}$, this is precisely \eqref{eq:negative-rising-shift}.
\end{proof}

\end{document}
